% ===================================================================
%  도표 제 1 호 — 학교, 고등학교, 교실.
%
%  Ceci n'est PAS une transcription : c'est une traduction, faite en
%  2026, en coreen d'aujourd'hui. D'ou trois differences avec
%  texto/io et texto/fr, et trois seulement :
%
%   * pas de \nl ni de \cc — la traduction ne suit aucune ligne
%     imprimee, et les lignes de ce fichier ne sont que du confort ;
%   * pas de \begin{VUpage} — elle n'a ni feuillet ni folio ;
%   * le gras se pose sur le TERME seul, ponctuation dehors.
%
%  Les cles « %%K » sont celles de texto/io, macro pour macro pour
%  l'apparat, et les renvois \textsuperscript{(n)} visent les MEMES
%  objets de la planche, dans le MEME ordre.
%
%  UNE COLONNE SANS PARENTE DANS TOUT LE RELEVE. Les trente-sept
%  colonnes precedentes se rangent en familles — romane, germanique,
%  slave, indo-aryenne, dravidienne, semitique, sinitique — et
%  chacune avait au moins une voisine. Le coreen n'en a aucune. Le
%  japonais lui ressemble par l'ordre des mots et par rien d'autre :
%  ni le lexique de base, ni la morphologie, ni l'ecriture. Le
%  chinois et le cantonais lui ont prete un lexique savant, et pas
%  une regle de grammaire.
%
%  CE QUI VAUT POUR LUI, EN REVANCHE, C'EST LA REGLE DU TELOUGOU. Le
%  coreen est a tete finale sans exception : tout ce qui qualifie
%  precede ce qui est qualifie, le verbe finit la phrase, et il n'y a
%  pas de proposition relative postposee. parigi.py sert donc ici
%  exactement comme il servait la — c'est la troisieme famille de
%  langues ou il sert, apres l'indo-aryen et le dravidien, et il n'a
%  pas fallu y toucher.
%
%  L'ECRITURE, ELLE, REND LE CONTROLE PLUS SIMPLE QU'AILLEURS. Le
%  cantonais devait se defendre contre le mandarin, le marathi contre
%  le hindi, l'egyptien contre l'arabe standard : trois voisins qui
%  partagent leur alphabet. Le hangul n'est partage par personne. Ce
%  qui guette cette colonne est autre chose — le mot sino-coreen
%  ecrit en ideogrammes, comme on l'imprimait encore en 1926. Le
%  coreen de 2026 s'ecrit tout en hangul ; un seul ideogramme dans
%  texto/ko est donc une faute, quel qu'il soit, et kolonoj.py la
%  releve d'une seule classe de caracteres.
%
%  L'APPARAT PREND SES PROPRES MOTS, comme pour les trois colonnes
%  precedentes : 도표 pour un tableau — 표 seul serait le mot le plus
%  banal de la langue —, 계열 pour une serie, et 마당 pour une scene.
%  마당 est le mot du theatre coreen, celui du 판소리 et du 마당놀이,
%  la ou le calque sino-coreen dirait 장 : troisieme colonne d'affilee
%  a preferer son propre mot de scene, apres le प्रवेश marathe et le
%  రంగం telougou.
%
%  LA CLASSE COREENNE A SON VOCABULAIRE ENTIER, et il est en hangul
%  jusqu'au bout :
%
%    칠판 (1)       le tableau noir       분필 (3)     la craie
%    칠판지우개 (2)  le torchon            석판 (42)    l'ardoise
%    잉크병 (25)    l'encrier             압지 (41)    le buvard
%    사전 (30)      le dictionnaire       공책 (28)    le cahier
%    연습장 (35)    le cahier de brouillon
%    자 (50)        la regle              컴퍼스 (26)  le compas
%    책가방 (57)    le cartable           필통 (31)    le plumier
%    책받침 (53)    le sous-main          수첩 (54)    le carnet
%    운동장 (75)    la cour               쉬는 시간 (76) la recreation
%    교장 (71)      le proviseur          학감 (72)    le censeur
%    수위 (77)      l'homme de service    장학사 (73)  l'inspecteur
%    들창 (78)      le vasistas
%
%  « 들창 » EST EXACTEMENT LE VASISTAS. Le mot est coreen de fond —
%  들다, soulever, et 창, fenetre : la petite fenetre haute qu'on
%  souleve pour l'air. Trente-huitieme colonne, trente-huitieme fois
%  que l'objet du renvoi 78 se laisse nommer sans periphrase.
%
%  « 수위 » EST LE MEME EMPLOI QUE LE فرّاش EGYPTIEN, LE शिपाई MARATHE
%  ET LE బంట్రోతు TELOUGOU : l'homme de service de l'ecole, celui qui
%  porte le registre des absents. Quatrieme ecole reelle d'affilee, et
%  la meme fonction nommee d'un mot.
%
%  LA GAMME RESTE OCCIDENTALE, ET C'EST VOULU. Le coreen a sa propre
%  notation — 궁상각치우 — et l'ecrire ici serait tentant. Mais le
%  tableau nomme do-re-mi et pose l'equivalence « = C. D. E. F. G. A.
%  B. » : c'est la gamme d'Europe qui est l'objet, et l'on ne
%  transpose pas les choses. Meme arbitrage qu'au telougou.
%
%  TROIS RETOURNEMENTS SUR LES HUIT PAIRES DE parigi.py : « stulo (8)
%  en sua katedro (6) », « nigra tabelo (1) kun la vishilo (2) » et
%  « furnelo (46) kun lua tubo (45) ». Les cinq autres sont des
%  enumerations ou des frontieres de phrase — les trois cas connus.
% ===================================================================

%%K t01-apar-0 apar
\VUsaut{2.0mm}
\VUcentre{12.6pt}{120}{설명 소책자}
\VUsaut{3.2mm}
\VUcentre{8.4pt}{160}{부속}
\VUsaut{2.6mm}
\VUtitre{15.6pt}{0}{델마 보조 도표}
\VUsaut{3.4mm}
\VUornamento{9pt}
\VUsaut{5.0mm}
\VUcentre{13.2pt}{90}{첫째 계열}
\VUsaut{3.0mm}
\VUfilet{16mm}
\VUsaut{5.6mm}

%%K t01-apar-1 apar
\VUtitre{15.0pt}{60}{도표 제 1 호}
\VUsaut{1.4mm}
\VUtitre{11.4pt}{0}{학교, 고등학교, 교실.}
\VUsaut{2.2mm}
\VUfilet{20mm}
\VUsaut{6.4mm}

%%K t01-tit-1 sub
\VUpk{10.2pt}{40}{전체 설명.}
\VUsaut{3.0mm}

%%K t01-01-1 p
1. --- 이 도표는 \VUgras{고등학교}의 \VUgras{교실}을 보여 준다. 이 교실에는
\VUgras{책상} \textsuperscript{(18)} 여덟 개와 \VUgras{걸상} \textsuperscript{(32)}
여덟 개가 있다. 걸상 하나에 학생 세 명이 앉는다. \VUgras{선생님} \textsuperscript{(7)}
께서는 \VUgras{의자} \textsuperscript{(8)}에 앉으시는데, 그 의자는 당신의
\VUgras{교단} \textsuperscript{(6)} 위에 있다.

%%K t01-02-1 p
2. --- 교단 왼쪽에는 \VUgras{유리장} \textsuperscript{(15)}이 \VUgras{끼워져} 있다.
오른쪽에는 \VUgras{칠판} \textsuperscript{(1)}이 있고, 거기에 \VUgras{칠판지우개}
\textsuperscript{(2)}와 \VUgras{분필} \textsuperscript{(3)}이 놓여 있다.

%%K t01-02-2 p
교단 위에는 \VUgras{벽걸이 그림표} \textsuperscript{(9, 11,} \textsuperscript{12)}와
\VUgras{지도} \textsuperscript{(10)}가 걸려 있다.

%%K t01-03-1 p
3. --- 학생이 모두 제 \VUgras{자리} \textsuperscript{(17)}에 앉아 있는 것은 아니다.
요안네스 \textsuperscript{(4)}는 칠판 앞에 서서 \VUgras{지우개}를 들고
\VUgras{기하 도형}을 지운다.

%%K t01-03-2 p
페트루스는 교단 가까이에 있다. 그는 \VUgras{쓰기 과제} \textsuperscript{(5)}를 걷어
선생님께 가져다 드린다.

%%K t01-04-1 p
4. --- \VUgras{이} 선생님은 \VUgras{현대어} 담당이시다. 그분은 학생들에게 외국어를
말하고 읽고 쓰는 법을 가르치신다 \VUgras{(독일어, 영어, 스페인어, 이탈리아어,
러시아어} 또는 \VUgras{프랑스어)}.

%%K t01-05-1 p
5. --- 교실은 아주 넓고 아주 밝다. 안쪽에는 \VUgras{들창} \textsuperscript{(78)} 두
개가 달린 넓은 \VUgras{창}과 \VUgras{유리문}이 있어서 바람을 통하게 하고 빛을 들인다.

%%K t01-05-2 p
이 교실은 춥지 않다. 한가운데에는 방을 데우려고 아주 좋은 \VUgras{난로}
\textsuperscript{(46)}가 있는데, 거기에 \VUgras{연통} \textsuperscript{(45)}이 달려
있다.

%%K t01-05-3 p
벽에는 \VUgras{옷걸이} \textsuperscript{(58)}가 박혀 있고, 거기에 학생들의 모자와
외투가 걸려 있다.

%%K t01-tit-2 sub
\VUsaut{5.2mm}
\VUpk{10.2pt}{40}{운동장}
\VUsaut{3.6mm}

%%K t01-06-1 p
6. --- \VUgras{시계} \textsuperscript{(63)}는 아홉 시 오 분을 가리킨다.

%%K t01-06-2 p
\VUgras{쉬는 시간} \textsuperscript{(76)}이 막 끝나고 수업이 다시 시작된다. 노는 곳은
\VUgras{운동장} \textsuperscript{(75)}이다. 놀고 있는 학생 몇이 보인다.
\VUgras{감독 선생} \textsuperscript{(74)}이 그들을 지켜본다.

%%K t01-07-1 p
7. --- 운동장에서는 여러 사람이 더 보인다. 곧 \VUgras{장학사} \textsuperscript{(73)},
학교장 \VUgras{(교장)} \textsuperscript{(71)}, 그리고 \VUgras{학감}
\textsuperscript{(72)}이다.

%%K t01-07-2 p
장학사는 학교를 살피고, 교장은 고등학교를 이끌고, 학감은 학업을 감독한다.

%%K t01-07-3 p
네 번째 사람은 \VUgras{수위} \textsuperscript{(77)}다. 그는 결석을 적어 넣으려고
겨드랑이에 장부를 끼고 다닌다.

%%K t01-08-1 p
8. --- 운동장 안쪽에는 \VUgras{체조 기구} \textsuperscript{(70)}와
\VUgras{수공 작업} \textsuperscript{(69)}을 위한 작업실이 있다.

%%K t01-08-2 p
도표 오른쪽으로는 \VUgras{음악}과 \VUgras{그리기} \VUgras{교실}이 보이기 시작한다.

%%K t01-noto-1 noto
\VUnotes{143.2mm}{%
(*) \textit{세례명}도 라틴어 꼴에 따라 옮겨 적기를 권한다. 그것만이 헤아릴 수 없는
갈래를 피하는 길이다. 게다가 \textit{Giovanni, Jean, Johann, Jan, Hans, John, Ivan}
따위 가운데에서 무엇을 고르겠는가? 세례명만을 위한 낱말집을 따로 만들 것인가? 라틴어에서
그 주격을 길어다가 이렇게 말하자 : \VUgras{Ioannes, Ioanna; Iulius, Iulia; Iakobus;}
\VUgras{Andreas; Lukas.} (Gram. Kompleta).%
}

%%K t01-tit-3 sub
\VUpk{10.2pt}{40}{자세한 설명.}
\VUsaut{2.4mm}

%%K t01-tit-4 sub
\VUpk{10.2pt}{40}{잘하는 학생들.}
\VUsaut{3.0mm}

%%K t01-09-1 p
9. --- 교단 오른쪽 첫째 걸상의 학생들은 \VUgras{헨리쿠스} (*), \VUgras{스테파누스},
\VUgras{카롤루스}다.

%%K t01-09-2 p
헨리쿠스는 아주 주의 깊고 부지런해서 선생님 말씀을 잘 듣는다.

자기 책상 위에 그는 \VUgras{공책} \textsuperscript{(28)}, \VUgras{책}
\textsuperscript{(29)}, \VUgras{사전} \textsuperscript{(30)}, \VUgras{필통}
\textsuperscript{(31)}을 두고 있다. 그의 책에는 \VUgras{쪽}이 많고, 장마다
\VUgras{문단}이 여럿이며, \VUgras{줄}마다 \VUgras{낱말}이 많다.

%%K t01-09-4 p
그 옆의 스테파누스 \textsuperscript{(24)}는 열심이다. 그는 다음 시간에 배울 것을 자기
공책에 적어 둔다. 그는 \VUgras{글씨체}가 곱다. 그의 책은 덮여 있어서 우리 눈에는
\VUgras{표지} \textsuperscript{(27)}만 보인다. 그 곁에는 그의 \VUgras{컴퍼스}
\textsuperscript{(26)}가 놓여 있다.

%%K t01-09-5 p
카롤루스 \textsuperscript{(23)}는 손으로 자기 책을 들고, 배운 것을 다시 읽으려고
그것을 편다. 그도 학생마다 그렇듯이 제 앞에 \VUgras{잉크병} \textsuperscript{(25)}을
두고 있다. 그는 말을 잘 듣는다.

%%K t01-10-1 p
10. --- 옆 책상에는 \VUgras{루도비쿠스}, \VUgras{안토니우스}, \VUgras{파울루스}가
있다. 루도비쿠스는 자기 \VUgras{배운 과} \textsuperscript{(22)}를 외우고 선생님의
\VUgras{물음}에 대답한다. 안토니우스 \textsuperscript{(21)}는 아주 주의가 없어서
이따금 선생님께 벌을 받기도 한다. 파울루스 \textsuperscript{(20)}는 좋은
\VUgras{동무}이고 상냥하며 남을 잘 돕는다. 그는 아주 주의 깊다.

%%K t01-tit-5 sub
\VUsaut{4.2mm}
\VUpk{10.2pt}{40}{못하는 학생들.}
\VUsaut{3.2mm}

%%K t01-11-1 p
11. --- 헨리쿠스 뒤에는 \VUgras{게오르기우스} \textsuperscript{(38)}가 있다. 나쁜
아이는 아니지만 \VUgras{수다스럽고} 주의가 없으며 장난이 심하다. 그의 \VUgras{경솔함}과
\VUgras{말 안 듣는 버릇}은 수업 중에 \VUgras{소란}을 일으켜서 자주 호된
\VUgras{꾸중}을 듣는다. 그의 \VUgras{연습장} \textsuperscript{(35)}은 지저분하다.
그는 자기 \VUgras{펜} \textsuperscript{(34)}으로 거기에 \VUgras{잉크를 묻히고}, 그래서
늘 \VUgras{지우개} \textsuperscript{(36)}를 쓴다. 그의 \VUgras{공책집}
\textsuperscript{(33)}은 찢어져 있는데, 그가 아주 조심성이 없기 때문이다. 그런데 왜
그는 현대어 교실에 자기 \VUgras{연필꽂이} \textsuperscript{(37)}를 들고 오는 걸까?

%%K t01-11-2 p
그 옆의 에드문두스 \textsuperscript{(39)}는 그를 좋아하지 않는다. 그는 사실 좀
굼뜨지만 말수가 적고 조용하다. 떠들지 않는다. 다만 듣는 대신에 자기
\VUgras{읽기 책}의 \VUgras{이야기}를 읽기를 더 좋아한다.

%%K t01-11-3 p
이 걸상의 셋째 학생은 \VUgras{루도비쿠스} \textsuperscript{(40)}다. 그는 부지런하다.
하얀 \VUgras{종이}에 글을 쓴다. 제 앞에는 \VUgras{압지} \textsuperscript{(41)}를
놓아 두었다.

%%K t01-12-1 p
12. --- 선생님 왼쪽 둘째 걸상의 학생들은 아주 어리다. 첫째인 어린 \VUgras{에밀리우스}는
아직 글씨를 잘 쓸 줄 몰라서 자기 \VUgras{석판} \textsuperscript{(42)}과
\VUgras{석필}과 \VUgras{스펀지} \textsuperscript{(43)}를 들고 다닌다.

%%K t01-12-2 p
그 곁에는 \VUgras{아우구스투스}와 \VUgras{베르나르두스} \textsuperscript{(44)}가
있다. 이 아이는 공부는 잘하지만 \VUgras{옷차림}이 아주 단정하지 못하다.

%%K t01-13-1 p
13. --- 그 뒤에는 \VUgras{게으른 아이} 셋이 있다. \VUgras{알베르투스}는 배운 것을
익히지 않는다. \VUgras{야코부스} \textsuperscript{(47)}는 듣는 대신에 잔다. 그의
\VUgras{게으름}은 \VUgras{꾸지람}과 \VUgras{벌}을 부른다. 끝으로
\VUgras{크리스티아누스} \textsuperscript{(48)}는 너무 진지하지 못하다. 그는 나쁘게
\VUgras{처신하고} 고치려고 조금도 애쓰지 않는다.

%%K t01-14-1 p
14. --- 그와 달리 그 이웃들은 행실이 바르다. \VUgras{에르네스투스}
\textsuperscript{(51)}는 질서와 규칙을 좋아해서 늘 자기 \VUgras{자}
\textsuperscript{(50)}와 \VUgras{연필} \textsuperscript{(49)}을 쓴다. 그는
\VUgras{통학생}이다.

%%K t01-14-2 p
\VUgras{프란치스쿠스} \textsuperscript{(52)}는 \VUgras{기숙생}이다. 제복을 입고
고등학교에서 먹고 잔다. 그는 좋은 \VUgras{동무}다.

%%K t01-14-3 p
마우리키우스는 자기 \VUgras{주머니칼} \textsuperscript{(56)}로 \VUgras{연필}을 깎는다.
그는 아주 꼼꼼하다.

%%K t01-14-4 p
그는 제 책을 모두 들고 다닌다. 자기 \VUgras{문법책} \textsuperscript{(55)}, 자기
\VUgras{수첩} \textsuperscript{(54)}, 심지어 \VUgras{책받침} \textsuperscript{(53)}
까지. 그의 \VUgras{책가방} \textsuperscript{(57)}은 뒤쪽 바닥에 놓여 있다.

%%K t01-14-5 p
교실 안쪽의 책상 두 개는 \VUgras{우등생} \textsuperscript{(19)}들이 쓴다.

%%K t01-tit-6 sub
\VUsaut{4.6mm}
\VUpk{10.2pt}{40}{그리기와 음악.}
\VUsaut{3.4mm}

%%K t01-15-1 p
15. --- \VUgras{미술실}에는 벽에 걸린 고운 \VUgras{본}들과, \VUgras{갓}이 달린
\VUgras{등} \textsuperscript{(62)}이 천장에 매달려 있다. \VUgras{선생님}
\textsuperscript{(60)}은 아주 참을성이 많으시다. 그분은 학생들의 \VUgras{그림}
\textsuperscript{(61)}을 고쳐 주시고, 실물을 보고 \VUgras{흉상} \textsuperscript{(59)}
그리는 법을 가르치신다.

%%K t01-16-1 p
16. --- \VUgras{음악실}은 아주 재미있어 보인다. 거기에서는 \VUgras{악보} 읽는 법,
\VUgras{오선} 위에 적힌 \VUgras{음표} \textsuperscript{(67)} 부르는 법(도, 레, 미,
파, 솔, 라, 시\,=\,C. D. E. F. G. A. B.), \VUgras{음자리표} 아는 법, 고운
\VUgras{가락} 부르는 법을 배운다. 악보는 \VUgras{보면대} \textsuperscript{(65)} 위에
놓는다. 학생 하나가 \VUgras{피아노를 치고} \textsuperscript{(64)},
\VUgras{음악 선생님} \textsuperscript{(66)}이 \VUgras{박자를 맞춘다}
\textsuperscript{(68)}.

%%K t01-17-1 p
17. --- \VUgras{수공 작업} \textsuperscript{(69)}을 위한 \VUgras{작업실}의 한 귀퉁이가
보이기 시작한다. 거기에서 아이들은 나무를 대패질하고 쇠를 줄질하는 법을 배울 수 있다.
모든 고등학교에 이런 곳이 있는 것은 아니다.

%%K t01-tit-7 sub
\VUsaut{5.0mm}
\VUpk{10.2pt}{40}{과학 교실.}
\VUsaut{3.6mm}

%%K t01-18-1 p
18. --- 수학 선생님은 학생들에게 \VUgras{기하}, \VUgras{산수}, \VUgras{셈},
\VUgras{미터법}을 가르치신다. 칠판 위에는 주된 기하 도형이 보인다 : \VUgras{원}
\textsuperscript{(a)}, \VUgras{정사각형} \textsuperscript{(b)}, \VUgras{삼각형}
\textsuperscript{(c)}, \VUgras{선} \textsuperscript{(d)}.

%%K t01-18-2 p
\VUgras{셈 하나}도 보인다. 그것은 \VUgras{덧셈}도 \VUgras{뺄셈}도 \VUgras{곱셈}도
아니고, 바로 \VUgras{나눗셈} \textsuperscript{(e)}이다.

%%K t01-19-1 p
19. --- 미터법 도표 위에는 이런 것이 보인다 : \VUgras{미터} \textsuperscript{(a)},
\VUgras{저울} \textsuperscript{(b)}과 그 \VUgras{분동}, \VUgras{리터}
\textsuperscript{(e)}, \VUgras{데카리터} \textsuperscript{(f)}, \VUgras{데시리터},
\VUgras{세제곱미터} \textsuperscript{(d)}.

%%K t01-19-2 p
학생들은 \VUgras{자연 과학} \textsuperscript{(11)}, 동물학 \textsuperscript{(a)},
식물학 \textsuperscript{(b)}, 지질학 \textsuperscript{(c)}, \VUgras{역사}
\textsuperscript{(12)}도 배운다.

%%K t01-19-3 p
유리장 안에는 \VUgras{물리} \textsuperscript{(15)}와 \VUgras{화학}
\textsuperscript{(16)} 기구가 있다.

%%K t01-19-4 p
유리장 위에는 \VUgras{구} \textsuperscript{(14)}, \VUgras{원뿔}, \VUgras{지구본}
\textsuperscript{(13)}이 있다.

%%K t01-19-5 p
도표들 위에는 세계의 다섯 대륙을 보여 주는 \VUgras{지도}가 걸려 있다 :
\VUgras{유럽} \textsuperscript{(g)}, \VUgras{아시아} \textsuperscript{(e)},
\VUgras{아프리카} \textsuperscript{(f)}, \VUgras{아메리카} \textsuperscript{(ab)},
\VUgras{오세아니아} \textsuperscript{(h)}, 그리고 두 큰 대양인 \VUgras{태평양}
\textsuperscript{(c)}과 \VUgras{대서양} \textsuperscript{(d)}.

\VUsaut{6.0mm}
\VUfilet{22mm}
